% Filtered-Only Evaluation of Market Regime Detection
% A Standardized Protocol and Cross-Method Benchmark Across Eight Historical Crises
%
% Categories: q-fin.CP primary, stat.ML secondary
% Target: arXiv preprint, ~15-20 pages single-column.

\documentclass[11pt]{article}

\usepackage[a4paper, margin=1in]{geometry}
\usepackage{amsmath, amssymb, amsthm}
\usepackage{graphicx}
\usepackage{booktabs}
\usepackage{hyperref}
\usepackage{natbib}
\usepackage{xcolor}
\usepackage{enumitem}

\hypersetup{
  colorlinks=true,
  linkcolor=black,
  citecolor=black,
  urlcolor=blue!50!black,
}

\title{%
  Filtered-Only Evaluation of Market Regime Detection: \\
  A Standardized Protocol and Cross-Method Benchmark \\
  Across Eight Historical Crises
}
\author{Kunal Jain \\ \texttt{kunal@spectatr.ai}}
\date{\today}

\begin{document}
\maketitle

\begin{abstract}
The literature on market-regime detection routinely reports detection-lag and
posterior-quality numbers using \emph{smoothed} posterior labels, which use
information from the entire observed sequence. This is appropriate for
descriptive analysis but inflates apparent performance for any real-time use
case. We propose a standardized \emph{filtered-only} evaluation protocol with
three components: (1) a static, cited registry of canonical historical crisis
events with multi-anchor dates (peak / $-5\%$ drawdown / $-10\%$ drawdown);
(2) a sustained-fire firing rule (3 consecutive days of $P(\text{crisis}) >
0.5$) to suppress flicker-credit; and (3) a calibration target rooted in
observable forward outcomes (21-day forward drawdown $> 10\%$) rather than
another model's regime estimate. We apply the protocol across six methods ---
Gaussian HMM, Markov-switching with Student-$t$ emissions, sparse jump model,
joint cross-sectional HMM with Fama-French factor regime-switching means,
Bayesian online change-point detection, and Wasserstein $k$-means --- on a
universe of US equity ETFs over [DATA RANGE]. We measure the gap between
filtered and smoothed detection lag at the level of individual historical
crises, and find [SUMMARY OF RESULTS — TO BE FILLED IN AT WEEK 12]. The
implementation is open-source and bit-deterministic under fixed seeds.
\end{abstract}

\section{Introduction}
\label{sec:intro}

% Hook: motivate filtered-vs-smoothed gap as the under-reported issue.
\textbf{TODO}: open with a single-paragraph framing of the smoothed-label problem.

% Contribution list.
Our contributions are:
\begin{itemize}[leftmargin=*]
  \item A standardized filtered-only evaluation protocol with multi-anchor
    crisis dates, sustained-fire firing rule, and observable-outcome
    calibration target.
  \item A cross-method benchmark applying the protocol to six regime-detection
    methods on US equity ETFs.
  \item An empirical measurement of the smoothed-vs-filtered detection-lag gap
    at the level of individual historical crises.
  \item An open-source bit-deterministic implementation
    (\url{https://github.com/<placeholder>}) covering the data layer, model
    family, ensemble, calibration, and backtest pipeline.
\end{itemize}

\section{Related Work}
\label{sec:related}

\textbf{TODO}: cite Hamilton (1989), Ang \& Bekaert (2002), Adams \& MacKay
(2007), Horvath et al.\ (2021), Liu, Maheu \& Song (2024), Shu (2024),
Frazzini, Israel \& Moskowitz (2018), Idzorek (2007), State Street (2025).

\section{Evaluation Protocol}
\label{sec:protocol}

\subsection{Crisis-Event Registry}
The benchmark uses eight historical events, anchored by static dates committed
in the implementation's \texttt{regime/eval/crises.py}: 2007 Quant Quake,
2008 Lehman / GFC, 2011 US Debt Ceiling, 2015 China Devaluation, 2018
Volmageddon, 2020 COVID, 2022 LDI / UK Gilt, 2024 Yen Carry Unwind. Anchor
dates are not derived from the underlying price data so the registry is
stable across data revisions.

\subsection{Detection-Lag Definition}
For a crisis with anchor $t^\star$, the detection lag of method $m$ is
\[
  L_m(t^\star) = \min \{t \geq t^\star : P_m(\text{crisis}\mid y_{1:t}) > \theta
                       \text{ for } \delta \text{ consecutive trading days}\} - t^\star,
\]
with default headline parameters $\theta = 0.5$ and $\delta = 3$
(Section~\ref{sec:protocol-anchors}). The headline anchor is the first close
at $-5\%$ from the rolling 90-day maximum (\texttt{m5\_date}); we report
multi-anchor sensitivity (peak / m5 / m10) and multi-threshold sensitivity
(0.3, 0.5, 0.7) in Appendix~\ref{appendix:sensitivity}.

\subsection{Calibration Target}
\label{sec:protocol-calibration}
The crisis head is calibrated against the observable indicator
\[
  y_t = \mathbb{1}\!\left[\max\{(P_t - P_{t+s}) / P_t : s = 1, \dots, 21\} > 0.10\right],
\]
where $P_t$ is SPY's adjusted close. Calibration uses isotonic regression on
$k$-fold out-of-fold logistic-regression outputs to avoid the in-sample
tautology; the trailing 21 days of each fold are masked because their
forward outcome is unobservable.

\subsection{Multi-Anchor Sensitivity}
\label{sec:protocol-anchors}
% TODO: table comparing detection lag under peak / m5 / m10 anchors with
% threshold = 0.5 and sustained-fire = 3.

\section{Methods Benchmarked}
\label{sec:methods}

We evaluate six methods. Four are state-based (output $P(s_t \mid y_{1:t})$);
two are change-point detectors (output native run-length / cluster-distance
features that are stacked into the calibrated crisis head, not stuffed into
a fake $K{=}3$ wrapper).

\paragraph{Gaussian HMM.} \textbf{TODO}: spec from \texttt{regime/models/hmm\_gaussian.py}.
\paragraph{Markov-switching Student-$t$ (MS-AR).} \textbf{TODO}: data-augmented EM.
\paragraph{Sparse Jump Model.} \textbf{TODO}: $k$-means with jump penalty + Viterbi DP.
\paragraph{Joint cross-sectional HMM.} \textbf{TODO}: factor-switching means $\alpha_k + B_k^{obs} f_t^{FF}$
+ rank-3 latent factor cov.
\paragraph{BOCPD.} \textbf{TODO}: Adams-MacKay 2007 with conjugate Normal-IG / Student-$t$ predictive.
\paragraph{Wasserstein $k$-means.} \textbf{TODO}: medoid clustering on rolling distributional windows.

\section{Data}
\label{sec:data}

\textbf{TODO}: universe of 18 US ETFs (SPY + 11 sector SPDRs + 6 style-factor
ETFs) + TLT, daily OHLCV from yfinance, Fama-French 5+Mom from Ken French
data library, FRED macro series. Description of point-in-time data layer
(action-log architecture, no retroactive mutation, snapshot-stable hashes
in \texttt{data.lock}).

\section{Results}
\label{sec:results}

\subsection{Detection-Lag Comparison}
\textbf{TODO}: Figure~\ref{fig:detection-lag} — per-crisis detection lag for
filtered vs smoothed across all six methods.

\subsection{Calibration}
\textbf{TODO}: Figure~\ref{fig:reliability} — reliability diagram for the
calibrated crisis head; PR-AUC with random baseline.

\subsection{Smoothed-vs-Filtered Gap}
\textbf{TODO}: Figure~\ref{fig:smoothed-vs-filtered} — the headline figure
showing the gap between $P(s_t \mid y_{1:t})$ and $P(s_t \mid y_{1:T})$
during a representative crisis.

\section{Failure Modes}
\label{sec:failure}

% Per project convention this section is longer than Section~\ref{sec:results}.
\textbf{TODO}: enumerate where the protocol disagrees with the published
literature; where individual methods break down; the period-2017 low-vol
grind regime where every method underperforms; the cost-sensitivity column
results.

\section{Reproducibility}
\label{sec:repro}

The implementation is bit-deterministic under fixed seeds and single-threaded
BLAS in CI. Every figure in this paper is regenerated by a script in
\texttt{paper/figures/}. To reproduce the entire paper from raw data:
\begin{verbatim}
git clone <repo>
uv sync --all-extras
just refresh        # fetch ETF + FRED + Fama-French data
just lock           # snapshot data hashes
make -C paper figures
make -C paper main.pdf
\end{verbatim}

\section{Conclusion}
\label{sec:conclusion}

\textbf{TODO}: one paragraph summarizing the protocol contribution and the
empirical finding about the filtered-vs-smoothed gap.

\bibliographystyle{plainnat}
\bibliography{refs}

\appendix

\section{Anchor / Threshold Sensitivity}
\label{appendix:sensitivity}
\textbf{TODO}: $3\times 3$ table of detection lag across (peak, m5, m10) anchors
and (0.3, 0.5, 0.7) thresholds.

\section{Pre-Registered Hyperparameters}
\label{appendix:hyperparameters}
The full set of hyperparameter values is committed in
\texttt{STRATEGY\_HYPERPARAMETERS.md} in the source tree; the pre-registration
commit is dated 2026-05-09 and precedes any backtest commit.

\end{document}
